\documentclass[a4paper,11pt]{article}
\usepackage[a4paper,margin=2cm]{geometry}
\usepackage[T1]{fontenc}
\usepackage[utf8]{inputenc}
\usepackage[ngerman,english]{babel}
\usepackage{lmodern,microtype,xcolor,parskip,enumitem}
\usepackage[most]{tcolorbox}
\usepackage{titlesec}
\usepackage[hidelinks]{hyperref}
\definecolor{HeaderBlueDark}{RGB}{70,110,170}
\definecolor{RowBlueLight}{RGB}{235,243,255}
\definecolor{RowBlueDark}{RGB}{220,233,250}
\definecolor{SoftGray}{RGB}{90,90,90}
\setlength{\parindent}{0pt}
\setlist[itemize]{leftmargin=1.4em,itemsep=0.35em,topsep=0.35em}
\titleformat{\section}[block]{\Large\bfseries\color{HeaderBlueDark}}{}{0pt}{}
\newtcolorbox{exbox}{enhanced,breakable,colback=RowBlueLight,colframe=RowBlueDark,boxrule=0.4pt,arc=1.5mm,left=1.2mm,right=1.2mm,top=1mm,bottom=1mm,title={Beispiele \textnormal{(Examples)}},fonttitle=\bfseries,colbacktitle=RowBlueDark,coltitle=black}
\NewDocumentEnvironment{grammarentry}{mm+b}{\begin{tcolorbox}[enhanced,breakable,colback=white,colframe=HeaderBlueDark,title={#1\hfill\normalfont\footnotesize #2},colbacktitle=HeaderBlueDark,coltitle=white,fonttitle=\bfseries]#3\end{tcolorbox}}{}
\newcommand{\GermanText}[1]{\textbf{Deutsch: }#1\par}
\newcommand{\EnglishText}[1]{{\small\color{SoftGray}\textbf{English: }#1\par}}
\newcommand{\ExampleItem}[2]{\item \textbf{#1}\\{\small\color{SoftGray}#2}}
\title{Die Modalpartikel \glqq schon\grqq\\{\large\textnormal{The modal particle schon}}}
\date{}
\begin{document}
\maketitle
\tableofcontents
\section{Bedeutung und Wirkung \textnormal{(Meaning and effect)}}
\begin{grammarentry}{schon}{Modalpartikel / Abtönungspartikel}
\GermanText{Schon kann beruhigen, Zuversicht ausdrücken oder in rhetorischen Fragen eine Möglichkeit herunterspielen.}
\EnglishText{Schon can reassure, express confidence or play down a possibility in rhetorical questions.}
\end{grammarentry}
\section{Grammatische Rolle und Satzarten \textnormal{(Grammar and sentence moods)}}
\begin{grammarentry}{Verwendung}{Satzart und Kontext}
\GermanText{Häufig in Zusicherungen und rhetorischen Fragen; die genaue Lesart hängt vom Tonfall ab.}
\EnglishText{Often in reassurances and rhetorical questions; the interpretation depends on intonation.}
\end{grammarentry}
\begin{grammarentry}{Form und Stellung}{Syntax}
\GermanText{Als Modalpartikel ist schon unflektierbar, fordert keinen Kasus und steht normalerweise im Mittelfeld. Allein im Vorfeld eines normalen Aussagesatzes kann eine Modalpartikel nicht stehen. Die Betonung und der Kontext beeinflussen die Wirkung.}
\EnglishText{As a modal particle, schon is invariable, does not govern case and normally appears in the middle field. It cannot occupy the prefield of a normal statement on its own. Stress and context affect its pragmatic force.}
\end{grammarentry}
\section{Typische Beispiele \textnormal{(Typical examples)}}
\begin{exbox}\begin{itemize}
\ExampleItem{Das wird schon klappen.}{It'll work out, don't worry.}
\ExampleItem{Du wirst schon sehen!}{You'll see!}
\ExampleItem{Was kann da schon passieren?}{What could possibly go wrong?}
\ExampleItem{Er wird schon kommen.}{I'm sure he'll turn up.}
\end{itemize}\end{exbox}
\section{Abgrenzung zu anderen Verwendungen \textnormal{(Contrast with other uses)}}
\begin{grammarentry}{Wortarten unterscheiden}{Nicht immer Modalpartikel}
\GermanText{Zeitliches schon bedeutet bereits und beschreibt, dass etwas früher als gedacht eingetreten ist.}
\EnglishText{Temporal schon means already and indicates that something has happened by now.}
\end{grammarentry}
\begin{exbox}\begin{itemize}
\ExampleItem{Der Zug ist schon da.}{The train is already here.}
\ExampleItem{Das wird schon!}{It'll be fine!}
\end{itemize}\end{exbox}
\section{Lernhinweis \textnormal{(Learning tip)}}
\GermanText{Eine Modalpartikel verändert vor allem die Gesprächsabsicht. Versuchen Sie nicht, jedem Vorkommen eine feste englische Eins-zu-eins-Übersetzung zu geben. Achten Sie auf Satzart, Tonfall und Situation.}
\EnglishText{A modal particle mainly affects the speaker's intention. Do not try to give it one fixed English translation in every context. Pay attention to sentence mood, intonation and situation.}
\section{Quellen \textnormal{(Sources)}}
\begin{itemize}
\item IDS Mannheim: \url{https://grammis.ids-mannheim.de/terminologie/2}
\item Duden: \url{https://www.duden.de/sprachwissen/sprachratgeber/Modalpartikeln}
\end{itemize}
\end{document}
